% Data request proposal for rail organizations.
% Generic by design: fill in the four fields below, and optionally the
% recipient, before sending. Build with:  latexmk -pdf data_request_proposal.tex
\documentclass[11pt]{article}

\usepackage[margin=1in]{geometry}
\usepackage[T1]{fontenc}
\usepackage{lmodern}
\usepackage{microtype}
\usepackage{booktabs}
\usepackage{tabularx}
\usepackage{array}
\usepackage{enumitem}
\usepackage{xcolor}
\usepackage[hidelinks]{hyperref}

% ---------------------------------------------------------------------------
% Fill these in before sending.
% ---------------------------------------------------------------------------
\newcommand{\authorname}{George James}
\newcommand{\authoraffil}{School of Engineering, Colorado State University Pueblo}
\newcommand{\contactemail}{\textcolor{red}{[your email]}}
\newcommand{\contactphone}{\textcolor{red}{[your phone]}}
\newcommand{\advisor}{\textcolor{red}{[advisor name and title]}}
\newcommand{\recipient}{your organization}   % e.g. "Example Railroad"
\newcommand{\usefulby}{\textcolor{red}{[date]}}
% ---------------------------------------------------------------------------

\setlist{itemsep=2pt, topsep=4pt}
\newcolumntype{L}{>{\raggedright\arraybackslash}X}
\setlength{\parskip}{4pt}
\setlength{\parindent}{0pt}

\title{\vspace{-1.5em}\textbf{Request for Freight Train Operating Data}\\[4pt]
  \large Validating a Longitudinal Train Dynamics Model for Fuel-Efficient Train Control}
\author{\authorname\\ \small\authoraffil}
\date{\today}

\begin{document}
\maketitle
\vspace{-1em}

\section*{Summary}

I am a graduate student building a fast, learned model of how freight trains
move: their speed, the forces between cars, and the energy used. The model is
meant to support fuel-efficient driving advice and planning. It is trained on
a physics-based simulator, and the most important open question in my work is
how closely that simulator matches real trains.

I am asking \recipient{} for a modest amount of recorded operating data
(locomotive event-recorder downloads and the matching train makeup) from a
handful of ordinary freight trips. I would use it only to check and calibrate
the model. I will handle it under whatever confidentiality terms you require,
and I will share the results with you. Section~\ref{sec:data} lists exactly
which data would help, in priority order. Even the first tier alone would be
valuable.

\section*{The research}

My master's thesis, \emph{A Graph Neural Surrogate for Longitudinal Train
Dynamics: Toward Fuel-Optimal Control of Freight Consists}, has three parts.
\begin{enumerate}
  \item \textbf{A reference simulator.} A longitudinal train dynamics model
    that treats every car as its own mass, joined by couplers with slack and
    draft gear. It includes traction and dynamic braking, rolling and curving
    resistance, and grade, on real route profiles built from public elevation
    and track data. It runs on a GPU, so thousands of trips can be simulated.
  \item \textbf{A learned surrogate.} A graph neural network that represents
    the train as a chain of cars and learns to reproduce the simulator. It
    is far faster than the simulator and handles any train length.
  \item \textbf{A control demonstration.} Using the surrogate to plan throttle
    and dynamic brake settings that reduce fuel use while respecting speed
    limits and coupler-force limits.
\end{enumerate}

\subsection*{Why real data matters}

Nearly every result in the thesis comes from the simulator. I have verified
that the code solves its equations correctly, and I am comparing it against
an international benchmark of published train dynamics simulators. Neither
shows that the equations describe a real train with real parameters. Measured
data is the only way to establish that, and it is what would let the fuel and
coupler-force results be stated with confidence rather than as results
``within the model.''

\section*{What data would help}
\label{sec:data}

The tiers are in priority order. Tier~1 by itself is enough to test the model
against reality. Anything partial is still useful, and I can work with
whatever format your systems export.

\subsection*{Tier 1: one trip, recorded}

For each trip, the locomotive event-recorder download (lead locomotive, and
remote units if the train runs distributed power) together with its consist
list.

\smallskip
\begin{tabularx}{\linewidth}{@{}>{\raggedright\arraybackslash}p{0.33\linewidth} l L@{}}
\toprule
\textbf{Parameter} & \textbf{Units} & \textbf{What it is used for} \\
\midrule
Time stamp & s & Aligning every channel \\
Speed & mph or km/h & The main quantity the model is checked against \\
Position (GPS or milepost) & lat/lon or mi & Locating the train on the grade and curve profile \\
Throttle notch & 0--8 & The driver's input; the model is run on exactly these commands \\
Dynamic brake setting & notch or \% & Same \\
Brake pipe pressure & psi & When and how hard the air brakes were applied \\
Brake cylinder pressure, if recorded & psi & Brake force actually developed \\
Independent (locomotive) brake & psi & Separating locomotive and train braking \\
Traction motor current or tractive effort, if recorded & A or lbf & Checking the locomotive force model directly \\
\bottomrule
\end{tabularx}

\medskip
\textbf{Consist list for the same trip}, front to rear:

\smallskip
\begin{tabularx}{\linewidth}{@{}>{\raggedright\arraybackslash}p{0.40\linewidth} L@{}}
\toprule
\textbf{Per vehicle} & \textbf{What it is used for} \\
\midrule
Position in train & The model is built car by car, in order \\
Locomotive model, or car type (e.g.\ covered hopper, tank) & Resistance, length and coupler type \\
Gross weight, or tare and load & Mass, the most important single parameter \\
Length over couplers & Spacing of masses along the train \\
Number of axles & Axle load, which enters rolling resistance \\
Distributed-power unit positions & Where remote locomotives sit \\
\bottomrule
\end{tabularx}

\medskip
\textbf{Sampling rate.} Whatever the recorder stores. One sample per second
is useful; faster is better.

\textbf{How much.} Five to twenty trips would be a strong data set. Variety
matters more than volume: loaded and empty trains, steep and level
territory, and if possible a trip or two with distributed power.

\subsection*{Tier 2: the route}

If the trips cover territory I cannot reconstruct from public data:
\begin{itemize}
  \item Track profile or track charts: grade (\%) and curvature (degrees or
    radius) by milepost.
  \item Timetable or posted speed restrictions over the same miles.
  \item A milepost-to-coordinate reference, if position is recorded only by
    milepost.
\end{itemize}

\subsection*{Tier 3: any of these, if they exist}

These are rarer, so I list them only in case they are on hand.
\begin{itemize}
  \item \textbf{End-of-train device data}: rear brake pipe pressure and
    motion, which show how a brake application travels along the train.
  \item \textbf{Fuel or energy}: fuel consumed per trip, or fuel rate from
    the locomotive. This is the quantity the control chapter tries to reduce.
  \item \textbf{Instrumented-car measurements}: coupler or drawbar forces
    from a load-cell car, at any position in a train.
  \item \textbf{Component data}: draft gear test curves, tractive-effort and
    dynamic-brake curves for your locomotive models, or car net braking ratios.
  \item \textbf{Resistance tests}: coast-down or other train resistance
    measurements.
\end{itemize}

\section*{How the data would be used}

\begin{itemize}
  \item \textbf{Validation.} Run the simulator on the recorded notch and brake
    commands over the recorded route, and compare its predicted speed with
    the measured speed along the whole trip.
  \item \textbf{Calibration.} Estimate the parameters that public sources give
    only roughly: train resistance, braking force, and locomotive force at
    each notch. Filtering methods (Kalman-type estimators) give these with
    statistical uncertainty, and test whether the model's errors behave like
    noise or show a systematic bias.
  \item \textbf{Reporting.} The thesis would report aggregate results only:
    error statistics and fitted parameters. No raw records or identifying
    details would be published.
\end{itemize}

\section*{Confidentiality and handling}

I understand operating data is sensitive. I will:
\begin{itemize}
  \item use the data only for this research, and not share or redistribute
    it;
  \item store it on encrypted storage that only I can access;
  \item accept data with train IDs, crew information, dates and car numbers
    removed. Locations can also be withheld if the grade and curvature
    profile along the trip is provided instead;
  \item work under a non-disclosure or data use agreement, arranged through
    the university if you prefer;
  \item give you the chance to review any passage that describes your data
    before the thesis is submitted, and acknowledge \recipient{} by name or
    not at all, as you choose;
  \item delete the data at the end of the project on request.
\end{itemize}

\section*{What you would receive}

\begin{itemize}
  \item A short report on how well the model reproduced your trains, with the
    calibrated parameters (resistance, braking and locomotive force
    estimates).
  \item A copy of the finished thesis.
  \item A briefing or call to walk through the results, if useful.
\end{itemize}

\section*{Timing and contact}

Data received by \usefulby{} can be included in the thesis. A single trip
with its consist list is enough to start, and I am happy to begin with a
small sample before anything larger is considered.

\medskip
\begin{tabular}{@{}l l@{}}
\authorname & \authoraffil \\
Email & \contactemail \\
Phone & \contactphone \\
Thesis advisor & \advisor \\
\end{tabular}

\medskip
Thank you for considering this request.

\end{document}
